\chapter{Banach 压缩映射原理}
\label{chap:banach-contraction}

许多方程没有可供直接计算的闭式解，却能改写成
\[
 x=T(x).
\]
若从一个初值 \(x_0\) 出发反复代入，
\(x_{n+1}=T(x_n)\)，便得到一条候选近似序列。要使这套办法成为定理，
至少要回答三个问题：迭代是否收敛，极限是否满足方程，方程是否只有一个
解。压缩映射原理用一个定量的距离缩短条件同时回答这三个问题；完备性则
保证逐步逼近不会收敛到空间外的缺口。

\section{不动点和迭代}

\begin{definition}[不动点与迭代轨道]
设 \(T:X\to X\)。满足 \(T(x)=x\) 的点称为 \(T\) 的不动点。给定
\(x_0\in X\)，递归定义
\[
 x_{n+1}=T(x_n)\qquad(n\ge0),
\]
所得序列称为从 \(x_0\) 出发的迭代轨道，也写作 \(x_n=T^n(x_0)\)。
\end{definition}

在方程求根中，写成 \(x=T(x)\) 的方式通常不止一种；不同的 \(T\)
可能有同样的不动点，却产生完全不同的迭代行为。例如方程 \(x^2=2\)
既可写成 \(x=x+(2-x^2)\)，也可写成
\[
 x=\frac12\left(x+\frac2x\right).
\]
前一种迭代在正根附近未必缩短误差，后一种则是常用的平方根迭代。

\begin{definition}[Lipschitz 映射与压缩映射]
设 \((X,d_X)\)、\((Y,d_Y)\) 为度量空间。若存在 \(L\ge0\)，使
\[
 d_Y(Tx,Ty)\le Ld_X(x,y)\qquad(x,y\in X),
\]
则称 \(T:X\to Y\) 为 Lipschitz 映射。若 \(X=Y\)，并且上式可取
\(0\le q<1\)，则称 \(T\) 为压缩映射，\(q\) 称为一个压缩常数。
\end{definition}

压缩常数不要求是所有可取常数中的最小值。条件中的同一个 \(q<1\) 必须
对所有点对成立；仅有
\[
 x\ne y\quad\Longrightarrow\quad d(Tx,Ty)<d(x,y)
\]
还不足以应用压缩映射原理。

Lipschitz 映射一致连续，因而连续。若 \(T\) 的 Lipschitz 常数为 \(L\)，
归纳可得
\[
 d(T^nx,T^ny)\le L^n d(x,y).
\]
特别地，压缩映射的高次迭代会以几何速度消除不同初值之间的距离。

\begin{proposition}[压缩映射至多有一个不动点]
\label{prop:contraction-fixed-point-unique}
压缩映射 \(T:X\to X\) 至多有一个不动点。
\end{proposition}

\begin{proof}
若 \(Tx=x\)、\(Ty=y\)，则
\[
 d(x,y)=d(Tx,Ty)\le qd(x,y).
\]
因为 \(1-q>0\)，只能有 \(d(x,y)=0\)，即 \(x=y\)。
\end{proof}

“至多一个”尚不说明存在性。映射 \(T(x)=x/2\) 在不完备空间
\((0,1)\) 上是压缩映射，却没有不动点；它的迭代趋于被空间删去的点
\(0\)。下一节的证明将指出，完备性恰好用于接住这条 Cauchy 轨道。

\section{压缩映射定理的完整证明}

\begin{theorem}[Banach 压缩映射原理]
\label{thm:banach-contraction}
设 \((X,d)\) 是非空完备度量空间，\(T:X\to X\) 是压缩映射，压缩
常数为 \(q<1\)。则 \(T\) 有唯一不动点 \(x^\ast\)。对任意初值
\(x_0\in X\)，迭代 \(x_{n+1}=T(x_n)\) 都收敛于 \(x^\ast\)。
\end{theorem}

\begin{proof}
任取 \(x_0\in X\)，令 \(x_{n+1}=T(x_n)\)。相邻项之间有
\[
 d(x_{n+1},x_n)
 =d(Tx_n,Tx_{n-1})
 \le q\,d(x_n,x_{n-1}),
\]
从而归纳得到
\begin{equation}
\label{eq:contraction-successive-step}
 d(x_{n+1},x_n)\le q^n d(x_1,x_0).
\end{equation}
若 \(m>n\)，三角不等式和几何级数给出
\begin{equation}
\label{eq:contraction-cauchy-tail}
\begin{aligned}
 d(x_m,x_n)
 &\le \sum_{k=n}^{m-1}d(x_{k+1},x_k)\\
 &\le \sum_{k=n}^{m-1}q^k d(x_1,x_0)
 \le \frac{q^n}{1-q}d(x_1,x_0).
\end{aligned}
\end{equation}
右端随 \(n\to\infty\) 趋于零，所以 \((x_n)\) 是 Cauchy 列。完备性
给出 \(x^\ast\in X\)，使 \(x_n\to x^\ast\)。

压缩条件还给出
\[
 d(Tx_n,Tx^\ast)\le qd(x_n,x^\ast)\longrightarrow0,
\]
即 \(Tx_n\to Tx^\ast\)。另一方面 \(Tx_n=x_{n+1}\to x^\ast\)。
极限的唯一性说明 \(Tx^\ast=x^\ast\)，所以不动点存在。唯一性由
命题~\ref{prop:contraction-fixed-point-unique} 得到。

上述构造从任意 \(x_0\) 开始都得到不动点，而不动点只有一个，故所有
迭代轨道都收敛于同一个 \(x^\ast\)。
\end{proof}

证明中各条件的作用可以逐行定位。压缩条件先控制相邻步长，再通过几何级数
使轨道成为 Cauchy 列；完备性只在“Cauchy 列在 \(X\) 中有极限”处
使用；最后再次使用压缩条件所含的连续性，把极限代入 \(T\)。因此这一定理
既是存在唯一性定理，也是带有明确算法的构造定理。

\begin{corollary}[闭不变子集上的压缩]
\label{cor:contraction-closed-invariant-subset}
设 \(X\) 完备，\(A\subset X\) 非空且闭。若 \(T:A\to A\) 是压缩
映射，则 \(T\) 在 \(A\) 中有唯一不动点，且从 \(A\) 中任意初值出发的
迭代都收敛到它。
\end{corollary}

\begin{proof}
闭子空间 \(A\) 完备，在 \(A\) 上应用
\theoremref{thm:banach-contraction} 即得结论。
\end{proof}

这个形式在应用中更常见。通常先找一个闭集合 \(A\)，证明 \(T(A)\subset A\)，
再在 \(A\) 上估计压缩常数。全空间之外的行为不影响结论。

\section{误差估计与完备性条件}

\subsection{先验估计、后验估计与残量}

压缩映射定理的证明已经包含误差界。令 \(x^\ast\) 为不动点。由
\eqref{eq:contraction-cauchy-tail} 令 \(m\to\infty\)，得到
\begin{equation}
\label{eq:contraction-apriori-error}
 d(x_n,x^\ast)
 \le \frac{q^n}{1-q}d(x_1,x_0).
\end{equation}
它只用初始步长和已知的 \(q\)，称为先验估计。若迭代已经进行到第
\(n\) 步，则从该步重新应用同一估计可得
\begin{equation}
\label{eq:contraction-aposteriori-error}
 d(x_n,x^\ast)
 \le \frac{1}{1-q}d(x_{n+1},x_n)
 \le \frac{q}{1-q}d(x_n,x_{n-1})\quad(n\ge1).
\end{equation}
这是后验估计：它用已经算出的相邻两项决定是否停止。

还可直接比较当前点与它的像。对任意 \(x\in X\)，
\[
\begin{aligned}
 d(x,x^\ast)
 &\le d(x,Tx)+d(Tx,Tx^\ast)\\
 &\le d(x,Tx)+q\,d(x,x^\ast),
\end{aligned}
\]
所以
\begin{equation}
\label{eq:contraction-residual-error}
 d(x,x^\ast)\le\frac{d(x,Tx)}{1-q}.
\end{equation}
量 \(d(x,Tx)\) 是不动点方程的残量。它小并不单独保证 \(x\) 接近某个
任意映射的不动点；分母 \(1-q\) 记录了压缩条件提供的稳定性。当
\(q\) 接近 \(1\) 时，小残量可能被明显放大。

若已知初值到不动点的距离，也有更直接的估计
\begin{equation}
\label{eq:contraction-direct-error}
 d(x_n,x^\ast)=d(T^nx_0,T^nx^\ast)
 \le q^n d(x_0,x^\ast).
\end{equation}
它形式最简，但 \(d(x_0,x^\ast)\) 往往正是未知量。

\subsection{两个不可删除的条件}

完备性不能省略。前述 \(T(x)=x/2\) 把 \((0,1)\) 映入自身，压缩常数
为 \(1/2\)，轨道却趋于 \(0\notin(0,1)\)。若工作集合不是全空间，要求
它闭也是为了保证它作为子空间完备。

统一的 \(q<1\) 同样不能换成逐点的严格缩短。令
\[
 X=[0,\infty),\qquad T(x)=x+\frac{1}{x+1}.
\]
空间 \(X\) 完备，且 \(T(x)>x\)，所以没有不动点。另一方面
\[
 T'(x)=1-\frac{1}{(x+1)^2}\in[0,1),
\]
由中值定理，对任意 \(x\ne y\) 都有
\[
 |T(x)-T(y)|<|x-y|.
\]
但 \(T'(x)\to1\)，不存在对所有点都适用的 \(q<1\)。这个例子说明，
证明所需的是可求和的几何衰减，而非每一步某种没有统一幅度的缩短。

\subsection{不动点对映射的稳定性}

\begin{theorem}[压缩映射的扰动估计]
\label{thm:contraction-perturbation}
设 \(X\) 完备，\(T,S:X\to X\) 均为压缩映射，其中 \(T\) 的一个压缩
常数为 \(q<1\)。设其不动点分别为 \(x_T,x_S\)。若存在
\(\eta\ge0\)，使
\[
 d(Tx,Sx)\le\eta\qquad(x\in X),
\]
则
\[
 d(x_T,x_S)\le\frac{\eta}{1-q}.
\]
\end{theorem}

\begin{proof}
利用两个不动点等式，并在 \(Tx_S\) 处插入中间项，
\[
\begin{aligned}
 d(x_T,x_S)
 &=d(Tx_T,Sx_S)\\
 &\le d(Tx_T,Tx_S)+d(Tx_S,Sx_S)\\
 &\le qd(x_T,x_S)+\eta.
\end{aligned}
\]
移项即得结论。
\end{proof}

这里不要求 \(S\) 与 \(T\) 具有同一个压缩常数；\(S\) 的压缩性只用来
保证 \(x_S\) 存在。若一族映射 \(T_\lambda\) 具有共同的压缩常数
\(q<1\)，且
\[
 d(T_\lambda x,T_\mu x)\le L\,\rho(\lambda,\mu)
 \qquad(x\in X),
\]
则相应不动点满足
\begin{equation}
\label{eq:parametric-fixed-point-bound}
 d(x_\lambda^\ast,x_\mu^\ast)
 \le\frac{L}{1-q}\rho(\lambda,\mu).
\end{equation}
这给出解对参数的 Lipschitz 连续依赖。

\section{方程求根、递推数列与含参数不动点}

\subsection{闭区间上的不动点迭代}

\begin{proposition}[一元压缩迭代]
\label{prop:scalar-contraction-iteration}
设 \(I=[a,b]\)，函数 \(\varphi:I\to I\) 可导，并满足
\[
 \sup_{x\in I}|\varphi'(x)|\le q<1.
\]
则方程 \(x=\varphi(x)\) 在 \(I\) 中有唯一解 \(\alpha\)，且对任意
\(x_0\in I\)，递推
\[
 x_{n+1}=\varphi(x_n)
\]
收敛于 \(\alpha\)。误差满足
\eqref{eq:contraction-apriori-error}--\eqref{eq:contraction-direct-error}
中相应估计。
\end{proposition}

\begin{proof}
中值定理给出
\[
 |\varphi(x)-\varphi(y)|\le q|x-y|.
\]
闭区间 \(I\) 完备，所以结论由压缩映射原理得到。
\end{proof}

以 \(\varphi(x)=\cos x\) 为例，它把 \([0,1]\) 映入
\([\cos1,1]\subset[0,1]\)，并且
\[
 |\varphi'(x)|=\sin x\le\sin1<1.
\]
因此 \(x=\cos x\) 在 \([0,1]\) 中有唯一解；从该区间任一点开始的递推
\(x_{n+1}=\cos x_n\) 都收敛。介值定理只给出解的存在，单调性可以给出
唯一性，而压缩映射原理还同时给出可执行的近似方法和误差界。

\subsection{单调方程的一种迭代格式}

\begin{theorem}[带导数界的求根迭代]
\label{thm:monotone-root-contraction}
设 \(F\in C^1([a,b])\)，并且存在常数 \(0<m\le M\)，使
\[
 m\le F'(x)\le M\qquad(a\le x\le b).
\]
若 \(F(a)\le0\le F(b)\)，则 \(F(x)=0\) 在 \([a,b]\) 中有唯一根
\(r\)。映射
\[
 T(x)=x-\frac{F(x)}{M}
\]
把 \([a,b]\) 映入自身，并且是压缩映射，压缩常数可取
\[
 q=1-\frac{m}{M}.
\]
因此任意 \(x_0\in[a,b]\) 所生成的 \(x_{n+1}=T(x_n)\) 都收敛于
\(r\)。
\end{theorem}

\begin{proof}
介值定理给出一个根；\(F'>0\) 说明 \(F\) 严格递增，所以根唯一。又
\[
 0\le T'(x)=1-\frac{F'(x)}{M}\le1-\frac{m}{M}=q<1,
\]
故 \(T\) 单调且满足压缩估计。

由中值定理，存在 \(\xi,\eta\in(a,b)\)，使
\[
 -F(a)=F'(\xi)(r-a)\le M(r-a),\qquad
 F(b)=F'(\eta)(b-r)\le M(b-r).
\]
于是
\[
 a\le T(a)\le r,\qquad r\le T(b)\le b.
\]
因 \(T\) 单调，
\[
 T([a,b])=[T(a),T(b)]\subset[a,b].
\]
现在可在完备空间 \([a,b]\) 上应用压缩映射原理。
\end{proof}

这个结果的导数下界 \(m>0\) 同时承担两项任务：保证根唯一，并把
\(T'\) 与 \(1\) 隔开。若只有 \(F'\ge0\)，则 \(q\) 可能等于 \(1\)，
几何误差估计随之消失。

\subsection{递推数列与线性迭代}

对递推数列，压缩映射原理把常见的“先证有界、再证单调、最后求极限”
换成另一条路线。它不要求序列单调；只要递推映射在一个完备不变集上压缩，
收敛和极限的唯一性便同时得到。

例如 \(a>0\)、\(x_0>\sqrt a\) 时，令
\[
 T(x)=\frac12\left(x+\frac{a}{x}\right),\qquad
 I=[\sqrt a,x_0].
\]
算术--几何平均不等式和 \(x\ge\sqrt a\) 给出
\[
 \sqrt a\le T(x)\le x,
\]
所以 \(T(I)\subset I\)。并且
\[
 0\le T'(x)=\frac12\left(1-\frac{a}{x^2}\right)
 \le\frac12\left(1-\frac{a}{x_0^2}\right)<1.
\]
故平方根迭代在 \(I\) 上压缩并收敛于唯一不动点 \(\sqrt a\)。这项证明
给出全程有效的几何误差界；若进一步利用迭代公式的特殊代数结构，还可得到
更快的二次误差估计。

在线性空间中，设 \(E\) 是有限维赋范空间，线性算子 \(A:E\to E\)
的算子范数
\[
 \|A\|:=\sup_{x\ne0}\frac{\|Ax\|}{\|x\|}
\]
满足 \(\|A\|\le q<1\)。对任意 \(b\in E\)，仿射映射
\[
 T(x)=Ax+b
\]
是压缩映射。因此 \(x=Ax+b\) 有唯一解，并且
\[
 x^\ast=\lim_{n\to\infty}
 \bigl(b+Ab+\cdots+A^{n-1}b+A^nx_0\bigr).
\]
取 \(x_0=0\) 得
\[
 x^\ast=\sum_{k=0}^{\infty}A^k b.
\]
这就是几何级数在算子情形的对应形式，也说明 \(I-A\) 可逆。

\subsection{含参数的不动点}

设参数空间为 \((\Lambda,\rho)\)，每个 \(\lambda\in\Lambda\) 对应
压缩映射 \(T_\lambda:X\to X\)。若所有映射共享同一个 \(q<1\)，
则 \eqref{eq:parametric-fixed-point-bound} 把映射随参数的变化直接传给
不动点。共同压缩常数不可由
“每个参数各有一个 \(q_\lambda<1\)”代替；若
\(q_\lambda\to1\)，估计中的 \(1/(1-q_\lambda)\) 会失去统一控制。

作为简单例子，若 \(A:E\to E\) 满足 \(\|A\|\le q<1\)，并令
\[
 T_b(x)=Ax+b,
\]
则不动点 \(x_b^\ast\) 对右端项满足
\[
 \|x_b^\ast-x_c^\ast\|
 \le\frac{1}{1-q}\|b-c\|.
\]
线性方程 \((I-A)x=b\) 的解因此不仅存在唯一，而且连续依赖于数据
\(b\)。以后讨论微分方程和积分方程时，同一估计会用于解对初值、参数和
右端项的稳定性。

\section*{习题}
\addcontentsline{toc}{section}{习题}

\subsection*{A组　压缩条件与不动点}
\begin{enumerate}[label=\textbf{\thechapter.\arabic*.}]
 \exerciseitem{contraction-continuity-and-uniqueness}
 证明 Lipschitz 映射一致连续；再直接证明压缩映射至多有一个不动点。
 \exerciseitem{iterate-contraction-factor}
 若 \(T\) 的压缩常数为 \(q\)，证明 \(T^n\) 的 Lipschitz 常数可取
 \(q^n\)。说明这并不意味着 \(q^n\) 必为最小 Lipschitz 常数。
 \exerciseitem{affine-contraction-neumann-series}
 设有限维赋范空间 \(E\) 上的线性算子满足 \(\|A\|<1\)。证明
 \(T(x)=Ax+b\) 有唯一不动点，并证明
 \[
  (I-A)^{-1}b=\sum_{k=0}^{\infty}A^kb.
 \]
 \exerciseitem{nonexpansive-fixed-point-failure}
 分别构造 Lipschitz 常数为 \(1\) 而有多个不动点、没有不动点的自映射，
 说明 \(q<1\) 不能放宽为 \(q\le1\)。
 \exerciseitem{closed-ball-invariance}
 设 \(T:X\to X\) 的压缩常数为 \(q<1\)。若
 \[
  d(Tx_0,x_0)\le(1-q)r,
 \]
 证明闭球 \(\overline B(x_0,r)\) 在 \(T\) 下不变。若 \(X\) 完备，
 由此得到球中的唯一不动点。
 \exerciseitem{closed-invariant-subset-contraction}
 设 \(X\) 完备，\(A\subset X\) 非空且闭，且 \(T(A)\subset A\)。
 若 \(T|_A\) 为压缩映射，证明它在 \(A\) 中有唯一不动点；说明
 \(T\) 在 \(X\setminus A\) 上是否压缩与结论无关。
\end{enumerate}

\subsection*{B组　证明与误差}
\begin{enumerate}[label=\textbf{\thechapter.\arabic*.},resume]
 \exerciseitem{banach-proof-reconstruction}
 不引用正文定理，从相邻项估计开始完整证明 Banach 压缩映射原理，并在
 证明中标出完备性使用的位置。
 \exerciseitem{contraction-apriori-aposteriori}
 推导先验估计 \eqref{eq:contraction-apriori-error} 和两个后验估计
 \eqref{eq:contraction-aposteriori-error}。若
 \(q=0.8\)、\(d(x_1,x_0)=1\)，求保证
 \(d(x_n,x^\ast)<10^{-6}\) 的一个先验步数。
 \exerciseitem{contraction-residual-bound}
 证明残量估计 \eqref{eq:contraction-residual-error}。构造 \(q\) 接近
 \(1\) 的仿射压缩映射，
 说明不能把系数 \(1/(1-q)\) 换成一个与 \(q\) 无关的常数。
 \exerciseitem{incomplete-space-contraction}
 在 \((0,1)\) 上分析 \(T(x)=x/2\)：证明它是压缩映射、没有不动点，
 并确定任意迭代轨道在 \(\R\) 中的极限。
 \exerciseitem{strict-distance-decrease-no-fixed-point}
 对 \(X=[0,\infty)\) 和 \(T(x)=x+1/(x+1)\)，证明
 \(x\ne y\) 时 \(|T(x)-T(y)|<|x-y|\)，但 \(T\) 没有不动点。
 说明为何不存在统一的压缩常数。
 \exerciseitem{contraction-perturbation-bound}
 完整证明扰动估计
 \(d(x_T,x_S)\le\eta/(1-q)\)，并给出等号成立的一维仿射例子。
 \exerciseitem{inexact-picard-iteration}
 设 \(T\) 的压缩常数为 \(q<1\)，\(x^\ast\) 为其不动点，而近似迭代满足
 \[
  d(x_{n+1},Tx_n)\le\varepsilon_n.
 \]
 证明
 \[
  d(x_n,x^\ast)
  \le q^n d(x_0,x^\ast)
    +\sum_{k=0}^{n-1}q^{\,n-1-k}\varepsilon_k.
 \]
 由此给出 \(\varepsilon_n\to0\) 时 \(x_n\to x^\ast\) 的证明。
 \exerciseitem{fixed-point-of-contractive-iterate}
 设 \(X\) 完备，\(T:X\to X\) 为映射，且某个 \(T^m\) 是压缩映射。
 证明 \(T\) 有唯一不动点，并且每条 \(T\)-迭代轨道都收敛到该点。
\end{enumerate}

\subsection*{C组　迭代模型}
\begin{enumerate}[label=\textbf{\thechapter.\arabic*.},resume]
 \exerciseitem{cosine-fixed-point-iteration}
 证明 \(\varphi(x)=\cos x\) 把 \([0,1]\) 映入自身且为压缩映射。
 从 \(x_0=0\) 出发，写出保证误差小于 \(10^{-6}\) 的先验步数表达式。
 \exerciseitem{monotone-root-fixed-point-scheme}
 在 \theoremref{thm:monotone-root-contraction} 的条件下，逐步证明
 \(T([a,b])\subset[a,b]\)，并推导根的迭代误差界。
 \exerciseitem{babylonian-square-root-iteration}
 对 \(a>0\)、\(x_0>\sqrt a\)，证明
 \[
  x_{n+1}=\frac12\left(x_n+\frac{a}{x_n}\right)
 \]
 在 \([\sqrt a,x_0]\) 上由压缩映射原理收敛于 \(\sqrt a\)。再由
 \[
  x_{n+1}-\sqrt a
  =\frac{(x_n-\sqrt a)^2}{2x_n}
 \]
 说明误差比几何估计所显示的还要快。
 \exerciseitem{parameter-dependent-fixed-points}
 设 \(T_\lambda:X\to X\) 共享压缩常数 \(q<1\)，且
 \[
  d(T_\lambda x,T_\mu x)\le L\rho(\lambda,\mu).
 \]
 证明不动点映射 \(\lambda\mapsto x_\lambda^\ast\) 是 Lipschitz 的。
 若右端只在 \(\mu\to\lambda\) 时一致趋于零，证明不动点映射连续。
 \optionalexerciseitem{optimal-affine-relaxation}
 设对称正定矩阵 \(A\) 的特征值都在 \([m,M]\subset(0,\infty)\)。
 对迭代
 \[
  x_{n+1}=x_n-\lambda(Ax_n-b)
 \]
 在欧氏范数下求使压缩常数
 \(\max_{\mu\in[m,M]}|1-\lambda\mu|\) 最小的 \(\lambda>0\)，并求该
 最小压缩常数。
 \exerciseitem{linear-data-stability}
 设 \(\|A\|\le q<1\)，\(x_b\) 是 \((I-A)x=b\) 的解。证明
 \[
  \|x_b-x_c\|\le\frac{\|b-c\|}{1-q},
 \]
 并说明该估计怎样由不动点扰动定理和 Neumann 级数分别得到。
\end{enumerate}
